% year: תשפ"ד
% semester: א
% moed: א
% number: 1
% points: 14
% categories: func-limits
% source: exams/מבחנים/תשפד/מועד א תשפד.pdf

\begin{question}
(14 נק') חשבו את הגבול הבא:
\[ \lim_{x\to\infty}\left(\frac{x^2+x+2}{x^2+x+1}\right)^{2x^2+4} \]
\end{question}

\begin{hint}
זהו גבול מהצורה $1^\infty$. כתבו את הבסיס בצורה $1+\frac{1}{x^2+x+1}$ והשתמשו בגבול $\lim_{t\to\infty}\left(1+\frac1t\right)^t=e$.
\end{hint}

\begin{solution}
כאשר $x\to\infty$ הבסיס שואף ל-$1$ והמעריך שואף ל-$\infty$, כלומר זהו ביטוי לא מוגדר מהצורה $1^\infty$.
נכתוב את הבסיס בצורה
\[ \frac{x^2+x+2}{x^2+x+1} = 1+\frac{1}{x^2+x+1}. \]
נסמן $t=x^2+x+1$; כאשר $x\to\infty$ גם $t\to\infty$. אז
\[ \left(1+\frac{1}{x^2+x+1}\right)^{2x^2+4} = \left[\left(1+\frac1t\right)^{t}\right]^{\frac{2x^2+4}{x^2+x+1}}. \]
הביטוי שבסוגריים המרובעים שואף ל-$e$ (הגבול המפורסם $\lim_{t\to\infty}(1+\frac1t)^t=e$), והמעריך מקיים
\[ \lim_{x\to\infty}\frac{2x^2+4}{x^2+x+1} = \lim_{x\to\infty}\frac{2+\frac{4}{x^2}}{1+\frac1x+\frac1{x^2}} = 2. \]
מרציפות הפונקציה $(u,v)\mapsto u^v=e^{v\ln u}$ עבור $u>0$ (או: לפי $e^{v\ln u}$ ורציפות $\ln$ ו-$\exp$) נקבל
\[ \lim_{x\to\infty}\left(\frac{x^2+x+2}{x^2+x+1}\right)^{2x^2+4} = e^{2}. \]

\textbf{דרך נוספת:} נכתוב את הביטוי כ-$e^{(2x^2+4)\ln\left(1+\frac{1}{x^2+x+1}\right)}$. מכיוון ש-$\ln(1+s)\sim s$ כאשר $s\to0$,
\[ (2x^2+4)\ln\left(1+\frac{1}{x^2+x+1}\right) \sim \frac{2x^2+4}{x^2+x+1}\xrightarrow[x\to\infty]{}2, \]
ושוב מרציפות האקספוננט מתקבל הגבול $\boxed{e^2}$.
\end{solution}
